% BOOK-FILE-ID: f-1debbece06cf
% section: 極限

\section{極限}\label{節：極限}
\sectionref*{節：ニュートン法},発展演習,\,\chapterref{章：解法の背景と広がり}で使う.基本パターン・応用パターンの一般項を求める段階では後回しにできる.

\subsection{極限の考え方}\label{節：極限の考え方}
$x$を$a$に限りなく近づけたとき,\,$f(x)$が近づく値を$\lim_{x\to a}f(x)$と書く.
例えば$\lim_{x\to2}(x^2-4)=0$である.

\subsection{数列の極限}\label{節：数列の極限}
番号$n$を大きくしたとき,\,$a_n$が$L$に限りなく近づくなら,\,$L$に収束するといい,
\[
  \lim_{n\to\infty}a_n=L
\]
と書く.例えば$\lim_{n\to\infty}\dfrac{1}{n}=0$である.

\subsection{関数の極限}\label{節：関数の極限}
$\lim_{x\to a}f(x)$では,\,$x=a$での値ではなく,\,$x\ne a$の近くでの値を見る.
極限の計算には,既知の極限を組み合わせる方法や,上下から挟む方法がある.
例えば$a_n\le b_n\le c_n$で,\,$a_n,c_n$が同じ$L$へ収束するなら,\,$b_n$も$L$へ収束する.
これをはさみうちの原理という.

\subsection{片側極限}\label{節：片側極限}
近づく側を限定したものを\textbf{\bookterm{one-sided}{片側極限}}という.
$x<a$から近づく\textbf{\bookterm{left-limit}{左側極限}}には
\[
  \begin{aligned}
    \lim_{x\to a-}f(x)&=\lim_{x\to a^-}f(x)\\
    &=\lim_{\substack{x\to a\\x<a}}f(x)=\lim_{x\uparrow a}f(x)
  \end{aligned}
\]
という記法がある.\,$x>a$から近づく\textbf{\bookterm{right-limit}{右側極限}}も同様に
\[
  \begin{aligned}
    \lim_{x\to a+}f(x)&=\lim_{x\to a^+}f(x)\\
    &=\lim_{\substack{x\to a\\x>a}}f(x)=\lim_{x\downarrow a}f(x).
  \end{aligned}
\]
$x\to0^+$を$x\to+0$,\,$x\to0^-$を$x\to-0$とも書く.\,$\pm0$は近づく側を表し,異なる数ではない.

左右から近づける場合,通常の極限が有限の$L$になるための必要十分条件は,左右の極限がともに存在し,\,$L$に一致することである.
例えば$x<0$で$f(x)=-1$,\,$x\ge0$で$f(x)=1$なら,左右の極限は$-1$と$1$で異なり,\,$\lim_{x\to0}f(x)$は存在しない.
片側極限にも,\,$x=a$での値そのものは関係しない.

\subsection{\texorpdfstring{ネイピア数$e$}{ネイピア数e}}\label{節：ネイピア数e}
ベルヌーイの定義では
\[
  e=\lim_{n\to\infty}\left(1+\frac1n\right)^n.
\]
\sectionref*{節：指数と対数}で用いた実数$t\to\infty$の表示でも同じ値になる.
以下の基本極限を,\bookterm{euler}{オイラーの公式}の導出に使う.

\subsection{三角関数の極限}\label{節：三角関数の極限}
角度を弧度法で表す.\,$0<x<\dfrac{\pi}{2}$では,単位円の扇形と内接・外接三角形の面積を比べると
\[
  \frac12\sin x<\frac12x<\frac12\tan x.
\]
\begin{center}
  \BookMathLayout{\begin{tikzpicture}[
        x=2.4cm,
        y=2.4cm,
        line cap=round,
        line join=round,
        >=stealth
      ]
      % x=40度を例に,扇形と内接・外接する三角形を描く.
      \pgfmathsetmacro{\figangle}{40}
      \coordinate (O) at (0,0);
      \coordinate (A) at (1,0);
      \coordinate (P) at ({cos(\figangle)},{sin(\figangle)});
      \coordinate (Q) at (1,{tan(\figangle)});

      % 第一象限だけを描き,各図形は境界線で示す.
      \draw[->,thin] (O)--(1.38,0) node[below] {$x$};
      \draw[->,thin] (O)--(0,1.28) node[left] {$y$};
      \draw[thin] (A) arc[start angle=0,end angle=90,radius=1];
      \draw[thick] (O)--(Q)--(A)--cycle;
      \draw[thick] (O)--(P)--(A);
      \draw[thick] (A) arc[start angle=0,end angle=\figangle,radius=1];
      \draw[->,semithick] (0.24,0)
      arc[start angle=0,end angle=\figangle,radius=0.24];

      \fill (O) circle (0.018);
      \fill (A) circle (0.018);
      \fill (P) circle (0.018);
      \fill (Q) circle (0.018);
      \node[below left] at (O) {$O$};
      \node[below] at (0.52,-0.03) {$1$};
      \node[below right] at (A) {$A$};
      \node[above left] at (P) {$P$};
      \node[above right] at (Q) {$Q$};
      \node at (0.55,0.20) {$\frac12\sin x$};
      \node at (0.27,0.07) {$x$};

      \node[anchor=east,align=right,text width=2.0cm] (sector-label)
      at (-0.10,0.88) {扇形\\[-1pt]面積 $\frac12x$};
      \draw[->,thin] (sector-label.east)--(0.94,0.342);
      \node[anchor=west,align=left,text width=2.1cm] (outer-label)
      at (1.33,0.62) {外接三角形\\[-1pt]面積 $\frac12\tan x$};
      \draw[->,thin] (outer-label.west)--(0.91,0.55);
    \end{tikzpicture}}{\resizebox{\linewidth}{!}{\begin{tikzpicture}[
          x=2.4cm,
          y=2.4cm,
          line cap=round,
          line join=round,
          >=stealth
        ]
        % x=40度を例に,扇形と内接・外接する三角形を描く.
        \pgfmathsetmacro{\figangle}{40}
        \coordinate (O) at (0,0);
        \coordinate (A) at (1,0);
        \coordinate (P) at ({cos(\figangle)},{sin(\figangle)});
        \coordinate (Q) at (1,{tan(\figangle)});

        % 第一象限だけを描き,各図形は境界線で示す.
        \draw[->,thin] (O)--(1.38,0) node[below] {$x$};
        \draw[->,thin] (O)--(0,1.28) node[left] {$y$};
        \draw[thin] (A) arc[start angle=0,end angle=90,radius=1];
        \draw[thick] (O)--(Q)--(A)--cycle;
        \draw[thick] (O)--(P)--(A);
        \draw[thick] (A) arc[start angle=0,end angle=\figangle,radius=1];
        \draw[->,semithick] (0.24,0)
        arc[start angle=0,end angle=\figangle,radius=0.24];

        \fill (O) circle (0.018);
        \fill (A) circle (0.018);
        \fill (P) circle (0.018);
        \fill (Q) circle (0.018);
        \node[below left] at (O) {$O$};
        \node[below] at (0.52,-0.03) {$1$};
        \node[below right] at (A) {$A$};
        \node[above left] at (P) {$P$};
        \node[above right] at (Q) {$Q$};
        \node at (0.55,0.20) {$\frac12\sin x$};
        \node at (0.27,0.07) {$x$};

        \node[anchor=east,align=right,text width=2.0cm] (sector-label)
        at (-0.10,0.88) {扇形\\[-1pt]面積 $\frac12x$};
        \draw[->,thin] (sector-label.east)--(0.94,0.342);
        \node[anchor=west,align=left,text width=2.1cm] (outer-label)
        at (1.33,0.62) {外接三角形\\[-1pt]面積 $\frac12\tan x$};
        \draw[->,thin] (outer-label.west)--(0.91,0.55);
      \end{tikzpicture}}}
  \BookFigureCaption{扇形と内接・外接三角形の面積比較}
\end{center}
整理すると$\cos x<\dfrac{\sin x}{x}<1$となる.
$x\to+0$では$\cos x\to1$なので,はさみうちの原理から右側極限は$1$である.
$\dfrac{\sin x}{x}$は偶関数なので左側も同じとなり,
\[
  \lim_{x\to0}\frac{\sin x}x=1.
\]
ここでは,扇形の面積公式と図形の包含関係を前提にしている.
$\dfrac{\tan x}{x}=\dfrac{\dfrac{\sin x}{x}}{\cos x}$より
\[
  \lim_{x\to0}\frac{\tan x}x=1
\]
も得られる.

\subsection{指数・対数の極限}\label{節：指数・対数の極限}
$\log x$は自然\bookterm{logarithm}{対数}$\log_e x$を表す.
まず$(1+x)^{\dfrac{1}{x}}$の極限を左右から求める.
$x\to0^+$では$t=\dfrac{1}{x}\to\infty$と置くと
\[
  (1+x)^{\dfrac{1}{x}}=\left(1+\frac1t\right)^t\longrightarrow e.
\]
$x\to0^-$では$t=-\dfrac{1}{x}\to\infty$,\,$u=t-1$と置き,
\begin{align*}
  (1+x)^{\dfrac{1}{x}}&=\left(1-\frac1t\right)^{-t}\\
  &=\left(1+\frac1u\right)^u
  \left(1+\frac1u\right)\longrightarrow e.
\end{align*}
対数の連続性と$\log e=1$から
\[
  \lim_{x\to0}\frac{\log(1+x)}x
  =\log e=1.
\]
また,実数$a\ne0$に対して,\,$n$が十分大きければ$1+\dfrac{a}{n}>0$であり,
\[
  n\log\left(1+\frac an\right)
  =a\frac{\log\left(1+\dfrac{a}{n}\right)}{\dfrac{a}{n}}\longrightarrow a.
\]
指数関数の連続性より
\[
  \lim_{n\to\infty}\left(1+\frac an\right)^n=e^a.
\]
$a=0$でも両辺は$1$である.
